%% 32장 여러 변수의 보간
\chapter{여러 변수의 보간}
\label{ch:multivariate-interp}

\abstract{변수가 둘 이상이면 보간의 모양이 자료의 배치에 따라 갈라진다. 자료가 격자 위에 있으면 한 변수의 방법을 축마다 차례로 쓰는 텐서곱으로 충분하고 그 계산은 크로네커 곱의 성질 덕분에 한 변수의 풀이를 되풀이하는 것으로 줄어든다. 자료가 흩어져 있으면 다항식 보간조차 유일하지 않을 수 있어 전혀 다른 도구가 필요하다. 이 장에서는 거리만의 함수를 기저로 쓰는 방사 기저 함수 보간을 양정치 함수와 조건부 양정치 함수의 이론으로 정당화하고 형태 매개변수를 고를 때 정확도와 조건수 사이에 피할 수 없는 맞바꿈이 생긴다는 것을 수치로 본다. 또 점들을 삼각형으로 나누어 조각마다 보간하는 선형 보간과 Clough--Tocher 보간을 보고, 격자와 흩어진 자료에서 수렴 차수를 잰다.}

\section{격자 위의 자료}
\label{sec:grid-interp}

\subsection{텐서곱}

자료가 축마다의 격자 $x_1<\dots<x_m$, $y_1<\dots<y_n$의 모든 점 $(x_i,y_j)$에서 값 $v_{ij}$로 주어지면, 한 변수의 보간을 축마다 차례로 쓰면 된다. 가장 단순한 \term{쌍선형 보간}{bilinear interpolation}은 $(x,y)$가 들어 있는 격자 칸의 네 꼭짓점 값으로
\[
p(x,y)=(1-t)(1-s)v_{00}+t(1-s)v_{10}+(1-t)s\,v_{01}+ts\,v_{11}
\]
를 계산한다. 여기서 $t$, $s$는 칸 안의 상대 위치이고 $v_{00}$ 등은 칸의 꼭짓점 값이다. 이 식은 먼저 $y$ 방향으로 두 번 선형 보간한 뒤 그 결과를 $x$ 방향으로 선형 보간한 것과 같고 순서를 바꾸어도 같다. 쌍선형 보간은 $x$와 $y$ 각각에 대해서는 일차식이지만 $xy$ 항이 있어 평면은 아니다. 변수가 $d$개이면 칸의 꼭짓점이 $2^d$개이고, 축을 하나씩 줄여 가며 계산하면 꼭짓점의 수가 반씩 줄어 일이 $O(2^d)$이다.

같은 생각으로 31장의 스플라인을 축마다 쓰면 \term{텐서곱 스플라인}{tensor-product spline} $s(x,y)=\sum_{i,j}c_{ij}B_i(x)B_j(y)$를 얻는다. 계수를 정하는 계도 축마다 나뉜다.

\begin{proposition}
\label{prop:tensor-collocation}
$x$축의 배치 행렬을 $A$, $y$축의 배치 행렬을 $B$라 하면 텐서곱 보간의 계수 행렬 $C$는 $ACB^{\mathsf T}=V$의 해이다. 따라서 $C$는 $V$의 각 열에 $A^{-1}$을 쓰고, 그 결과의 각 행에 $B^{-1}$을 써서 얻는다.
\end{proposition}

\begin{proof}
보간 조건은 $\sum_{k,l}c_{kl}B_k(x_i)B_l(y_j)=v_{ij}$, 곧 $\sum_{k,l}A_{ik}c_{kl}B_{jl}=v_{ij}$이므로 $ACB^{\mathsf T}=V$이다. $A$와 $B$가 정칙이면 $C=A^{-1}V(B^{\mathsf T})^{-1}$이다. 벡터화하면 이것은 계 $(B\otimes A)\operatorname{vec}C=\operatorname{vec}V$이고, 크로네커 곱의 역은 각 인수의 역의 크로네커 곱이다.
\end{proof}

곧 $mn\times mn$ 계를 직접 풀 필요가 없다. 31장의 띠 풀이를 $x$축의 줄 $n$개와 $y$축의 줄 $m$개에 한 번씩 쓰면 되고 줄끼리는 독립이므로 병렬로 할 수 있다. 수렴 차수도 한 변수의 것을 물려받는다. $[0,1]^2$에서 시험 함수로 흔히 쓰는 Franke 함수\cite{franke1982}를 $n\times n$ 칸의 격자로 보간하면, 내부 $[0.1,0.9]^2$에서 잰 최대 오차가 $n=10,20,40,80$에서 쌍선형은 $6.8\times10^{-2}$, $1.9\times10^{-2}$, $4.9\times10^{-3}$, $1.2\times10^{-3}$으로 $O(h^2)$, 쌍3차 스플라인은 $2.0\times10^{-2}$, $7.2\times10^{-4}$, $3.5\times10^{-5}$, $2.0\times10^{-6}$으로 $O(h^4)$이다.

\subsection{차원의 저주}

텐서곱의 약점은 차원이다. 축마다 $n$개의 점을 두면 격자점이 $n^d$개이다. $d=10$에서 축마다 $10$개만 두어도 $10^{10}$개의 값이 필요하다. 높은 차원의 함수를 근사하는 일은 격자로는 불가능하고 함수의 구조를 쓰는 다른 방법이 필요하다. 흩어진 자료의 방법들은 점의 개수가 차원과 무관하다는 점에서 이 문제를 조금 비켜 간다.

\section{흩어진 자료와 방사 기저 함수}
\label{sec:rbf}

\subsection{다항식의 한계}

자료점 $x_1,\dots,x_N\in\mathbb R^d$이 아무렇게나 흩어져 있으면 다항식 보간은 문제가 많다. $d\ge2$에서는 점의 개수가 다항식 공간의 차원과 맞아도 보간이 유일하지 않을 수 있다. 예컨대 평면의 점 여섯 개가 한 원 위에 있으면, 그 원의 방정식이 이차 다항식이므로 보간하는 이차 다항식에 그것을 더해도 여전히 보간한다. 한 변수에서 Vandermonde 행렬이 언제나 정칙이었던 것과 달리, 여러 변수에서는 고정된 기저로는 모든 점 배치에서 정칙인 보간을 만들 수 없다는 것이 알려져 있다(Mairhuber--Curtis 정리\cite{wendland2004}). 그러므로 기저 자체가 자료점에 따라 달라져야 한다.

\subsection{방사 기저 함수}

\term{방사 기저 함수}{radial basis function, RBF} 보간은 자료점마다 그 점까지의 거리만의 함수를 하나씩 두고
\[
s(x)=\sum_{j=1}^Nc_j\,\phi(\|x-x_j\|)
\]
로 쓴다. 보간 조건 $s(x_i)=y_i$는 $A_{ij}=\phi(\|x_i-x_j\|)$에 대한 계 $Ac=y$이다. Hardy가 지형도를 그리려고 다중 이차식 $\sqrt{1+r^2}$을 쓴 것이 시작이다\cite{hardy1971}. $A$가 정칙인지는 $\phi$의 성질로 정해진다. 함수 $\Phi:\mathbb R^d\to\mathbb R$가 서로 다른 모든 점과 $0$이 아닌 모든 $c$에 대해
\[
\sum_{j,k}c_jc_k\Phi(x_j-x_k)>0
\]
이면 \term{양정치 함수}{positive definite function}라 한다. 그러면 $A$는 양정치 행렬이고 보간은 유일하게 존재한다.

\begin{proposition}
\label{prop:fourier-pd}
연속이고 적분 가능한 $\Phi$의 푸리에 변환 $\hat\Phi$가 음이 아니고 $0$이 아닌 연속함수이면 $\Phi$는 양정치이다. 특히 가우스 함수 $e^{-\varepsilon^2\|x\|^2}$는 모든 차원에서 양정치이다.
\end{proposition}

\begin{proof}
푸리에 역변환으로 $\Phi(x)=(2\pi)^{-d}\int\hat\Phi(\omega)e^{i\omega\cdot x}\,d\omega$로 쓰면
\[
\sum_{j,k}c_jc_k\Phi(x_j-x_k)=(2\pi)^{-d}\int\hat\Phi(\omega)\Bigl|\sum_jc_je^{i\omega\cdot x_j}\Bigr|^2d\omega\ge0
\]
이다. 등호가 성립하면 $\hat\Phi>0$인 열린 집합 위에서 $\sum_jc_je^{i\omega\cdot x_j}=0$이다. 이 함수는 $\omega$의 해석 함수이므로 모든 $\omega$에서 $0$이고, 서로 다른 점의 지수함수 $e^{i\omega\cdot x_j}$는 일차독립이므로 $c=0$이다. 가우스 함수의 푸리에 변환은 다시 가우스 함수이므로 양수이다.
\end{proof}

역방향, 곧 양정치 함수의 푸리에 변환이 비음 측도라는 것은 Bochner의 정리이다\cite{bochner1933}. 거리만의 함수가 모든 차원에서 양정치일 조건은 Schoenberg가 완전 단조 함수로 특징지었다\cite{schoenberg1938}. 가우스 함수, 역다중이차식 $1/\sqrt{1+r^2}$, 역이차식 $1/(1+r^2)$이 이런 함수이다.

\subsection{조건부 양정치 함수}

그러나 실제로 많이 쓰는 함수들, 예컨대 $-r$, $r^3$, $r^2\log r$은 양정치가 아니다. 이들은 다항식을 더해 고친다. 차수 $m-1$ 이하의 다항식에 대해 $\sum_jc_jp(x_j)=0$인 모든 $0$이 아닌 $c$에서 이차 형식이 양수이면 $\Phi$를 차수 $m$의 \term{조건부 양정치}{conditionally positive definite}라 한다. 그런 함수에는 차수 $m-1$의 다항식 항을 붙여
\[
s(x)=\sum_jc_j\phi(\|x-x_j\|)+\sum_kb_kp_k(x),\qquad\sum_jc_jp_k(x_j)=0
\]
으로 쓰고 덧붙인 조건 덕분에 미지수와 식의 개수가 맞는다. 행렬로 쓰면
\begin{equation}
\label{eq:rbf-saddle}
\begin{pmatrix}A&P\\P^{\mathsf T}&0\end{pmatrix}\begin{pmatrix}c\\b\end{pmatrix}=\begin{pmatrix}y\\0\end{pmatrix},\qquad P_{jk}=p_k(x_j)
\end{equation}
이다. Micchelli는 $-r$이 차수 $1$, 박판 스플라인 $r^2\log r$과 $r^3$이 차수 $2$, $-r^5$가 차수 $3$의 조건부 양정치임을 보였다\cite{micchelli1986}.

\begin{proposition}
\label{prop:rbf-saddle}
$\Phi$가 차수 $m$의 조건부 양정치이고, 자료점에서 차수 $m-1$ 이하의 다항식이 $0$이면 그 다항식이 $0$뿐이라고 하자(곧 $P$의 열이 일차독립이다). 그러면 \eqref{eq:rbf-saddle}의 행렬은 정칙이다.
\end{proposition}

\begin{proof}
$Ac+Pb=0$, $P^{\mathsf T}c=0$이라 하자. 첫 식에 $c^{\mathsf T}$를 곱하면 $c^{\mathsf T}Ac+(P^{\mathsf T}c)^{\mathsf T}b=c^{\mathsf T}Ac=0$이다. $P^{\mathsf T}c=0$이므로 조건부 양정치성에 의해 $c=0$이다. 그러면 $Pb=0$이고, $P$의 열이 일차독립이므로 $b=0$이다.
\end{proof}

이 행렬은 대칭이지만 부정이다. 그래서 Cholesky 대신 20장의 Bunch--Kaufman 대칭 부정 분해로 푼다. 조건부 양정치성은 행렬 $A$ 자체의 고유값에서도 드러난다. 차수 $m$의 조건부 양정치 $\Phi$에서 $A$의 이차 형식은 여차원 $q=\dim\mathcal P_{m-1}$인 부분공간에서 양수이므로, 4장의 Courant--Fischer 정리에 의해 $A$의 음의 고유값은 $q$개 이하이다. 평면의 점 $200$개에서 박판 스플라인의 $A$를 만들어 보면 실제로 양의 고유값이 $197$개, 음의 고유값이 $3$개이다. $q=\dim\mathcal P_1=3$과 정확히 맞는다.

박판 스플라인에는 이름 그대로 31장의 \cref{thm:holladay}를 평면으로 옮긴 성질이 있다. 보간하는 모든 함수 가운데 굽힘 에너지 $\iint(f_{xx}^2+2f_{xy}^2+f_{yy}^2)$를 최소로 하는 것이 박판 스플라인이다\cite{duchon1977}. 휘어지는 얇은 금속판을 점들에 고정했을 때의 모양이다.

\subsection{형태 매개변수와 불확정성 원리}

가우스 함수처럼 척도가 있는 핵에서는 $\phi(\varepsilon r)$의 \term{형태 매개변수}{shape parameter} $\varepsilon$을 골라야 한다. $\varepsilon$이 크면 기저 함수가 뾰족해져 각자 자기 점 근처만 보고, 작으면 평평해져 서로 닮는다. 평평한 기저는 매끄러운 함수를 아주 정확하게 근사할 수 있지만 기저 함수들이 거의 같아지므로 행렬 $A$의 조건수가 폭발한다. Schaback은 이것이 피할 수 없는 맞바꿈이어서 오차 상계와 조건수의 곱이 아래로 유계라는 \term{불확정성 원리}{uncertainty principle}를 보였다\cite{schaback1995}.

평면의 Halton 점\cite{halton1960} $200$개에서 Franke 함수를 가우스 RBF로 보간하면 \cref{tab:rbf-eps}와 같다.

\begin{table}[t]
\begin{sdtable}{colspec={ccc}, row{1}={font=\small\bfseries}}
\toprule
$\varepsilon$ & $\kappa_2(A)$ & 최대 오차 \\
\midrule
$12$ & $3.5\times10^{2}$ & $4.6\times10^{-1}$ \\
$8$ & $9.4\times10^{4}$ & $1.9\times10^{-1}$ \\
$6$ & $7.7\times10^{7}$ & $5.8\times10^{-2}$ \\
$4$ & $4.4\times10^{13}$ & $7.3\times10^{-1}$ \\
$3$ & $9.8\times10^{17}$ & $7.9$ \\
$2$ & $5.3\times10^{19}$ & $23$ \\
\bottomrule
\end{sdtable}
\caption{가우스 RBF 보간의 형태 매개변수와 조건수, $[0,1]^2$의 $81\times81$ 격자에서 잰 최대 오차. Halton 점 $200$개, numpy의 LU로 풀었다.}
\label{tab:rbf-eps}
\end{table}

$\varepsilon$을 $12$에서 $6$으로 줄이면 오차가 여덟 배 줄지만 조건수는 $10^8$에 이르고, 더 줄이면 반올림이 해를 삼켜 오차가 오히려 커진다. 정확한 산술이라면 $\varepsilon\to0$의 극한은 다항식 보간으로 가는 매끄러운 근사이지만 binary64에서는 그 정보를 잃는다. 같은 자료에 일차 다항식을 붙인 박판 스플라인을 쓰면 형태 매개변수 없이 오차가 $2.6\times10^{-2}$, 조건수가 $4.2\times10^4$이다. 척도가 없는 핵의 장점이다. Fornberg와 Wright는 $\varepsilon$이 작은 영역을 복소 $\varepsilon$의 등고선 적분으로 안정하게 계산하는 방법을 제안했다\cite{fornberg2004}.

전역 RBF는 $N\times N$의 밀집 계를 풀어야 하므로 $N$이 수만이면 감당하기 어렵다. 실제로는 평가점마다 가장 가까운 자료점 $k$개만으로 작은 RBF를 만드는 국소 판을 쓴다. 가까운 점을 찾는 일은 70장의 KD 트리가 맡는다. 대신 이웃 집합이 바뀌는 곳에서 보간이 연속하지 않다.

\section{삼각분할 위의 보간}
\label{sec:triangulation-interp}

\subsection{조각 선형}

평면의 흩어진 점들을 꼭짓점으로 하는 삼각형들로 볼록 껍질을 나누면, 삼각형마다 세 꼭짓점의 값을 잇는 평면으로 보간할 수 있다. 점 $x$가 꼭짓점 $a,b,c$의 삼각형 안에 있으면 $x=\lambda_aa+\lambda_bb+\lambda_cc$, $\lambda_a+\lambda_b+\lambda_c=1$인 \term{무게중심 좌표}{barycentric coordinates} $\lambda\ge0$이 하나로 정해지고, 보간값은 $\lambda_af(a)+\lambda_bf(b)+\lambda_cf(c)$이다. 무게가 비음이므로 결과는 세 값의 최솟값과 최댓값 사이에 있다. 이 보간은 연속이지만 삼각형의 변에서 꺾이고 오차는 삼각형의 크기 $h$에 대해 $O(h^2)$이다.

같은 점들의 삼각분할은 여럿이고 어느 것을 고르느냐에 따라 보간의 질이 다르다. 가늘고 긴 삼각형에서는 멀리 떨어진 점끼리 값을 잇게 되기 때문이다. 표준은 어떤 삼각형의 외접원 안에도 다른 점이 없는 \term{Delaunay 삼각분할}{Delaunay triangulation}으로, 가능한 모든 삼각분할 가운데 가장 작은 각을 최대로 한다\cite{sibson1978}. 그 계산과 정확한 기하 술어는 71장에서 다룬다.

\subsection{Clough--Tocher 보간}

변에서 꺾이지 않는 $C^1$ 보간을 원하면 삼각형 하나에 다항식 하나로는 부족하다. Clough와 Tocher는 유한요소법을 위해 각 삼각형을 무게중심에서 세 개의 작은 삼각형으로 나누고 작은 삼각형마다 3차 다항식을 두는 방법을 만들었다\cite{clough1965}. 꼭짓점의 값과 기울기, 그리고 각 변의 가운데에서 잰 법선 방향 도함수를 맞추면, 조각들이 삼각형 안에서 $C^1$으로 이어지고 이웃 삼각형과도 변을 따라 $C^1$으로 이어진다.

자료에는 기울기가 없으므로 먼저 꼭짓점마다 기울기를 추정해야 한다. scipy는 Nielson과 Renka--Cline의 방법을 따라, 삼각분할의 변을 따라 곡선들을 정의하고 그 굽힘 에너지가 최소가 되는 기울기를 반복법으로 구한다\cite{nielson1983,renka1984}. 이렇게 하면 같은 Franke 함수를 Halton 점으로 보간할 때 내부 $[0.1,0.9]^2$의 최대 오차가 점 $100$, $400$, $1600$, $6400$개에서 조각 선형은 $9.2\times10^{-2}$, $2.5\times10^{-2}$, $7.2\times10^{-3}$, $2.7\times10^{-3}$이고, Clough--Tocher는 $3.7\times10^{-2}$, $7.6\times10^{-3}$, $1.0\times10^{-3}$, $3.3\times10^{-4}$이다. 점이 네 배가 되어 간격이 반이 될 때마다 선형은 대략 네 배씩, Clough--Tocher는 대략 일곱 배씩 준다.

\subsection{가장 가까운 점}

가장 단순한 흩어진 자료의 보간은 가장 가까운 자료점의 값을 쓰는 것이다. 결과는 자료점들의 Voronoi 영역마다 상수인 계단 함수이다. 거리의 역수를 무게로 한 평균을 쓰는 Shepard의 방법은 연속이지만 자료점에서 기울기가 $0$이 되는 평평한 자리가 생긴다\cite{shepard1968}. 두 방법 모두 근사로서는 거칠지만 범주형 자료나 매끄러움을 가정할 수 없는 자료에서 쓰인다.

\section{stardust의 다변수 보간}
\label{sec:multivariate-stardust}

\begin{instardust}[\code{docs/design/interpolate.md} §2.7, §3.7]
stardust의 \code{RegularGridInterpolator}의 선형 보간은 꼭짓점 무게의 곱을 더하는 scipy와 달리, 마지막 축부터 첫 축까지 한 변수의 선형 보간을 차례로 쓴다. 수학적으로는 같지만 이 꼴은 노드에서 정확하고 결과가 꼭짓점 값의 최솟값과 최댓값 사이에 있으며 1차원에서 \code{sd.interp}와 같은 비트이다. \code{cubic}과 \code{quintic}은 \cref{prop:tensor-collocation}대로 축마다 31장의 띠 풀이를 써서 계수를 직접 구한다. scipy는 같은 계를 반복법으로 풀어서 설계 문서의 기록에 따르면 이차식 자료의 $6\times6$ 격자에서 참값 $25$ 대신 $24.999999999999986$을 낸다.
\end{instardust}

\begin{instardust}[\code{docs/design/interpolate.md} §2.8, §3.8]
\code{RBFInterpolator}는 scipy의 여덟 핵을 같은 식으로 계산하고, \eqref{eq:rbf-saddle}의 계를 대칭으로 채워 LAPACK의 \code{sytrf}(Bunch--Kaufman)로 분해한다. 다항식 항은 자료를 상자 $[-1,1]^d$로 옮긴 좌표의 단항식이다. 풀기 전에 $P$의 랭크를 SVD로 확인해 \cref{prop:rbf-saddle}의 가정이 깨지면 에러를 낸다. scipy는 이 경우, 예컨대 직선 위의 점들에서 일차 다항식 항을 쓰면 말없이 $10^{11}$ 크기의 값을 낸다. 같은 자료점이 두 번 있으면 행렬의 두 행이 같아 특이하므로, 평활 매개변수가 $0$이면 에러를 낸다. 삼각분할 위의 보간은 71장에서 볼 stardust 자체의 정확한 Delaunay 삼각분할 위에서 돈다.
\end{instardust}

\section{맺음}

여러 변수의 보간에서는 자료의 배치가 방법을 정한다. 격자 위에서는 한 변수의 이론이 축마다 그대로 통하지만 점의 개수가 차원에 대해 지수적으로 는다. 흩어진 자료에서는 고정된 기저가 실패하고 자료점에 붙은 방사 기저 함수나 자료점이 만드는 삼각분할이 그 자리를 채운다. RBF의 이론은 양정치 함수라는 해석학의 개념에 기대고 실제 계산에서는 형태 매개변수에 따라 커지는 조건수가 늘 문제가 된다. 다음 장에서는 보간의 다른 쪽 끝, 곧 함수의 값으로 도함수를 어림하는 수치 미분을 본다.

흩어진 자료의 근사 이론은 Wendland의 책\cite{wendland2004}이 엄밀하고, 실제 계산은 Fasshauer의 책\cite{fasshauer2007}이 친절하다.

\chapterreferences
